\chapter{Introduction}
\label{ch:introduction}

\section{Motivation}
% Single-core performance of HPC kernels depends less on what a loop computes than on the
% shape in which it reaches the compiler: source language, generator (translator, DaCe,
% LLM agent) and compiler flags. HPCAgent-Bench grades every representation of a kernel
% elementwise against one NumPy reference, which fixes *what* is computed and makes the
% representation the only free variable.
% Hook: same kernel, same GCC backend, 6--12x between representations (LLR-40 headline).

\section{Research Questions}
% RQ1 (bullet 1a) How do numba, C, C++ and Fortran lowerings of the same kernel compare
%     on loop-level-reasoning kernels (LLR-40), against hand-written C and agent code?
% RQ2 (bullet 1b) Do those findings transfer to kernels extracted from real applications
%     (CoMet, QuaTrEx, SpBench, WarpX, GraphAIBench, ANMLZoo), which are irregular,
%     mostly non-float64, and not flat loop nests?
% RQ3 (bullet 2) Under elementwise grading, which solver optimizations are legal at all,
%     and can a skill steer an agent to them (preconditioner freedom verdict)?
% RQ4 (bullet 3) What remained open in the DaCe codegen analysis of the CloudSC loop
%     nests, and what do the open items show once completed?

\section{Contributions}
% - LLR-40 measurement matrix: 40 kernels x 6 representations, pinned protocol, all
%   individual timings retained, 229/240 cells validated.
% - Eight application kernels extracted into HPCAgent-Bench across six Berkeley dwarfs
%   (PRs #11, #17, #18, #20, #21, #22), each with NumPy and native references.
% - Language comparison extended to the extracted kernels.
% - Solver-preconditioners skill (3 versions) with measured evidence and evaluations,
%   plus 14 solver test kernels it was developed against.
% - CloudSC loop-nest analysis completed/extended beyond Bonsall's semester project.

\section{Thesis Outline}
% One line per chapter, with \Cref once the labels are filled in.
